% from ega4-2-fr.tex
\oldpage[IV-2]{162}

\begin{proposition}[6.11.8]
\label{IV.2.s6.11.8-fr}
Soit $X$ un préschéma localement noethérien vérifiant la propriété suivante:
\begin{description}
\item[(CMU)] Tout sous-préschéma fermé intègre $Y$ de $X$ contient un ouvert
non vide $W$ tel que le préschéma induit par $Y$ sur $W$ soit un préschéma
de Cohen--Macaulay.
\end{description}
Alors, pour tout $\mathcal O_X$-Module cohérent $\mathcal F$, la fonction
$x\mapsto\operatorname{coprof}(\mathcal F_x)$ est localement constructible
et semi-continue supérieurement; les ensembles $U_{C_n}(\mathcal F)$ et
$U_{S_n}(\mathcal F)$ sont ouverts dans $X$.
\end{proposition}

\begin{proof}
En effet, soit $Y$ un sous-préschéma fermé intègre de $X$, de point
générique $y$; en vertu de (6.10.6), il y a dans $Y$ un voisinage ouvert
$V$ de $y$ tel que, pour tout $x\in V\cap Y$, on ait
\begin{equation}
\operatorname{coprof}(\mathcal F_x)=
\operatorname{coprof}(\mathcal F_y)+
\operatorname{coprof}(\mathcal O_{Y,x}).
\tag{6.11.8.1}\label{IV.2.s6.11.8.1-fr}
\end{equation}
Mais par hypothèse il existe un ouvert non vide $W$ de $Y$ tel que, pour
$x\in V\cap W$, on ait $\operatorname{coprof}(\mathcal O_{Y,x})=0$, donc
$\operatorname{coprof}(\mathcal F_x)$ est constante dans un voisinage de
$y$ dans $Y$, ce qui prouve que la fonction
$x\mapsto\operatorname{coprof}(\mathcal F_x)$ est localement constructible
(0$_{\rm III}$, 9.3.2); en outre il résulte alors de (6.11.5) et de
(0$_{\rm III}$, 9.3.4) que cette fonction est semi-continue supérieurement.
La dernière assertion résulte de (6.11.4, (i)), ou encore de (6.11.6).
\end{proof}

\begin{remarks}[6.11.9]
\label{IV.2.s6.11.9-fr}
\begin{enumerate}
\item[(i)] Si $X$ vérifie l'hypothèse (CMU) de (6.11.8), alors, pour tout
morphisme $u:X'\to X$ localement de type fini, $X'$ vérifie aussi (CMU).
Soient en effet $Y'$ un sous-préschéma fermé intègre de $X'$, $y'$ son point
générique, $y=u(y')$, et soit $Y$ le sous-préschéma intègre de $X$ ayant
pour espace sous-jacent $\overline{\{y\}}$; alors $u|Y'$ se factorise en
$Y'\xrightarrow{v}Y\xrightarrow{j}X$, où $j$ est l'injection canonique
(I, 5.2.2), et $v$ est localement de type fini (I, 6.6.6). En se
restreignant à des voisinages ouverts affines de $y$ et $y'$ respectivement,
on peut donc se borner au cas où $X=\operatorname{Spec}(A)$, où $A$ est un
anneau intègre de Cohen--Macaulay, et $X'=\operatorname{Spec}(A')$, où $A'$
est un anneau intègre contenant $A$ et qui est une $A$-algèbre de type fini.

